\section{Atributos de calidad}
\label{sec:atributos_de_calidad}

    \subsection{Identificación de los QA claves}
    \label{sec:identificacion_de_los_qa_claves}

        Se identifican los siguientes atributos de calidad como claves en el contexto del negocio.
        Adicionalmente, se citan especificaciones y requerimientos que motivaron estas afirmaciones.

        \subsubsection{Availability (disponibilidad)}
        \label{sec:disponibilidad}

            Habilidad del sistema para reparar o enmascarar fallas de manera tal que el período fuera de servicio no exceda un determinado valor durante un lapso de tiempo especificado.

            \begin{customenv}[frametitle=]
                La startup arVault es una fintech que opera una billetera digital
            \end{customenv}

            Al ser una billetera digital, el servicio permite manipular el dinero de los usuarios. Un servicio con baja o nula disponibilidad implicaria que los usuarios no pueden acceder a su propio dinero.
            Adicionalmente, el modelo de negocio se basa en que los usuarios operen, si el servicio no es disponible para permitir estas operaciones, se estan perdiendo potenciales ganancias.

        \subsubsection{Performance}
        \label{sec:performance}

            Habilidad del sistema para reaccionar ante ciertos eventos en un determinado tiempo.

            \begin{customenv}[frametitle=]
                Si bien el lanzamiento fue exitoso, comenzaron a aparecer reclamos de los usuarios sobre problemas en el uso de la función de cambio de monedas.
            \end{customenv}

            Uno de los principales problemas que impulsaron esta review es el endpoint de /exchanges, como se mostro es (TODO) este considerablemente mas lento que el resto de endpoints, lo que podria implicar el descontento de los usuarios.

            \begin{customenv}[frametitle=]
                las tasas de cambio no tienen gap entre el valor de compra y de venta.
            \end{customenv}

            Adicinalmente, el modelo de negocios es ofrecer una tasa de cambio sin gap, es decir en tiempo real.
            Una performance pobre en esta operacion implicaria que las tasas podrian que tener que ajustarse negativamente en comparacion a lo que se le ofrece al usuario en primera instancia.

        \subsubsection{Scalability (escalabilidad)}
        \label{sec:escalabilidad}

            Habilidad del sistema para incorporar recursos adicionales de manera efectiva. Efectividad = mejora mensurable de algún QA + no requerir un esfuerzo inapropiado + no interrumpir el funcionamiento.

            \begin{customenv}[frametitle=]
                Para ganar usuarios, las tasas de cambio...
            \end{customenv}

            Se especifica que objetivo a corto plaso es incrementar el numero de usuarios, el cual, el sistema debe ser capaz de manejar incorporando los recursos necesarios.

        \subsubsection{Reliability (confiabilidad)}
        \label{sec:confiabilidad}

            Grado en que la arquitectura es susceptible a fallas a nivel de sistema cuando se producen fallas parciales en componentes, conectores o datos.

            \todo{TODO: justificar}

        \subsubsection{Security (seguridad)}
        \label{sec:seguridad}

            Capacidad del sistema de prevenir acciones maliciosas o accidentales fuera del uso esperado.

            \begin{customenv}[frametitle=]
                se valida que haya saldo en la cuenta del cliente, se espera que lo haga la UI
            \end{customenv}

            La posibilidad de que los usuarios puedan realizar acciones maliciosas implican perdidas monetarias para el resto de usuarios, para la propia organizacion, ademas de una perdida de reputacion y con esto perdida de usuarios e inversiones.

    \subsection{Otros QA menos relevantes}
    \label{sec:otros_qa_menos_relevantes}

        Ademas se mencionan otros QA, que aunque no resulten claves y por ende no son una prioridad, son desceables, especialmente conforme el proyecto crezca.

        \todo{TODO: busca en el enunciado o en la documentacion como justificar estos QA}

        \subsubsection{Efficiency}
        \label{sec:eficiencia}

        Performance por recurso consumido.

        \subsubsection{Testability (testeabilidad)}
        \label{sec:testeabilidad}

        Grado en que un componente/sistema de software permite ser probado para exhibir sus fallas; usualmente relacionado al esfuerzo necesario para llevar adelante las tareas de testing.

        \subsubsection{Manageability}
        \label{sec:manejabilidad}

        Define la facilidad de gestión del sistema por parte de los administradores del mismo.

        \subsubsection{Simplicity}
        \label{sec:simplicidad}

        Facilidad con la que la arquitectura de un sistema puede ser entendida, razonada y explicada.

    \subsection{Decisiones de diseño actuales y su impacto en los QA}
    \label{sec:decisiones_impacto_qa}

        \subsubsection{Impacto en la Performance}
        \label{sec:impacto_en_la_performance}

        ...

    \subsection{Critica}
    \label{sec:critica}
